% Chapter 2 - BASIC.

\guidechapter{PROGRAMMING IN BASIC}{%
  \item Installing the wedge
  \item The twenty-three keywords
  \item Files and directories
  \item The RAM expansion
  \item Turbo
  \item Talking to the Ultimate directly
  \item What BASIC cannot do
}

\section{Installing The Wedge}

The wedge comes two ways, and they are the same program.

\textbf{As a program.} \bas{load"uci",8} then \bas{run}. It installs itself
into the 4K block at \cmd{\$C000}, which BASIC does not use, and leaves your
38911 bytes free exactly as they were.

\textbf{As a cartridge.} Mount \cmd{uci.crt} from the Ultimate's menu and the
wedge is there the moment the machine starts, with nothing to load. The
cartridge costs BASIC its 8K: you get 30719 bytes free instead of 38911.

Either way the machine greets you and the keywords are live. They survive
\bas{new}, they survive \bas{run/stop}~+~\key{restore}, and they survive a
program crashing. They do not survive a reset.

\begin{note}
The wedge takes over four of BASIC's own vectors --- the ones for tokenising,
listing, statement dispatch and expression evaluation. Every other wedge and
fast-loader that does the same thing will fight with it. If you use one, load
the SDK's wedge last.
\end{note}

\section{The Twenty-Three Keywords}

They fall into four groups, and you can ignore three of them until you need
them.

\subsection*{Asking how things went}

Nothing in this wedge stops your program with an error. A command that fails
sets a result code and returns, because a demo dropping to \bas{ready.} in the
middle of a part is worse than one that quietly did nothing. So every command
is followed, when it matters, by a look at \bas{uerr}.

\begin{reftable}{lp{0.62\linewidth}}
\bas{uerr} & the result of the last command. \bas{0} is success \\
\bas{udev} & the raw code the Ultimate reported, for diagnosis \\
\bas{ulen} & how many bytes came back \\
\bas{ust\$} & the status text, exactly as the Ultimate sent it \\
\bas{udat\$} & the reply as a string, up to 255 bytes \\
\bas{ubyte(n)} & byte \bas{n} of the reply, or \bas{0} past the end \\
\end{reftable}

\subsection*{Files, and the RAM expansion}

\begin{reftable}{lp{0.62\linewidth}}
\bas{udir} & print the current directory \\
\bas{uload} \textit{name\$} [,\textit{addr}] & load a PRG; no address means the
  file's own \\
\bas{ubload} \textit{name\$}, \textit{addr}, \textit{len} & load raw bytes, and
  not one past \textit{len} \\
\bas{usave} \textit{name\$}, \textit{addr}, \textit{len} & memory to a file \\
\bas{ustash} \textit{addr}, \textit{reu}, \textit{len} & C64 memory into the
  expansion \\
\bas{ufetch} \textit{addr}, \textit{reu}, \textit{len} & and back out of it \\
\bas{ureu} & the expansion's size in 64K banks, \bas{0} for none \\
\end{reftable}

\subsection*{Speed}

\begin{reftable}{lp{0.62\linewidth}}
\bas{uturbo} \textit{n} & set the processor speed index \\
\bas{uturbo} & as a function, the current index --- or \bas{255} when the
  machine's owner has turned turbo off \\
\end{reftable}

\subsection*{Everything else}

\begin{reftable}{lp{0.62\linewidth}}
\bas{uci} \textit{t}, \textit{c} [, \textit{args}] & send any command to any
  target \\
\bas{uci} & on its own, rescue a wedged interface \\
\bas{uw(}\textit{x}\bas{)} & force a 16-bit argument \\
\bas{ul(}\textit{x}\bas{)} & force a 32-bit argument \\
\bas{udos1} \bas{udos2} \bas{unet} \bas{uctrl} \bas{uiec} \bas{uhttp} &
  the six target numbers, by name \\
\end{reftable}

\section{Files And Directories}

\bas{udir} prints the directory and is the quickest way to see that everything
is working:

\begin{screen}
UDIR
\end{screen}

Loading is the interesting one, because the Ultimate has a fast path and the
SDK takes it when it can. This loads a file to the address stored in the file
itself, exactly as \bas{load"x",8,1} would:

\begin{screen}
10 ULOAD "GAME"\\
20 IF UERR THEN PRINT "COULD NOT LOAD": END
\end{screen}

Give it an address and it loads there instead:

\begin{screen}
10 ULOAD "SPRITES",832\\
20 IF UERR THEN PRINT "NO SPRITES": END
\end{screen}

\bas{ubload} is the one to use for data rather than programs. It interprets
nothing, strips no header, and stops at the length you give it:

\begin{screen}
10 UBLOAD "MUSIC.BIN",49152,4096
\end{screen}

\begin{note}
Filenames go to the Ultimate exactly as you type them. Uppercase PETSCII and
uppercase ASCII are the same bytes, so a name typed at a C64 arrives as the
Ultimate expects it, and the filesystem does not care about case anyway.
\end{note}

\section{The RAM Expansion}

If there is an expansion fitted, you have up to sixteen megabytes to put things
in --- and moving bytes in and out of it costs almost nothing, because the REU
does it with the processor halted rather than a byte at a time.

Always ask first. \bas{ureu} returns the size in 64K banks and zero when there
is nothing there, so one test covers both questions:

\begin{screen}
10 B = UREU\\
20 IF B = 0 THEN PRINT "NO EXPANSION": END\\
30 PRINT "EXPANSION IS"; B * 64; "K"
\end{screen}

Now put something in it and get it back. This copies the top of the screen into
the expansion, clears the screen, and puts it back:

\begin{screen}
10 IF UREU = 0 THEN PRINT "NEED AN REU": END\\
20 USTASH 1024,0,1000\\
30 PRINT CHR\$(147)\\
40 FOR I = 1 TO 2000: NEXT\\
50 UFETCH 1024,0,1000
\end{screen}

Line~20 sends 1000 bytes from screen memory at 1024 to expansion address~0.
Line~50 brings them back. Both are instant.

\begin{note}
The expansion address is a full 24-bit number, so it can reach all sixteen
megabytes. BASIC handles that fine --- \bas{ustash 1024, 1048576, 256} is a
perfectly good line.
\end{note}

\begin{warning}
Writing past the end of the expansion does not fail --- it wraps round to the
beginning and quietly overwrites what is there. That is exactly how \bas{ureu}
measures the size, and it is why asking first is worth the one line.
\end{warning}

\section{Turbo}

On an Ultimate~64 the processor can run faster than 1MHz, and \bas{uturbo}
sets it. The catch is that this is the machine owner's decision, not your
program's: unless \bas{turbo control} is set in the Ultimate's menu, the
registers are not there at all.

\begin{screen}
10 UTURBO 3\\
20 IF UERR THEN PRINT "NO TURBO HERE"\\
30 PRINT "SPEED INDEX IS"; UTURBO
\end{screen}

Index \bas{0} is 1MHz and index \bas{3} is 4MHz on every machine that has
turbo. Above~3 the numbers mean different speeds on different machines, so a
program that ships to other people should stay at or below~3 or ask.

The function form answers \bas{255} when turbo is unavailable. That is not a
speed --- the index is only four bits --- so you can test for it.

\section{Talking To The Ultimate Directly}

The wedge wraps the operations most programs want. Everything else --- and
there are around a hundred commands --- is reachable with \bas{uci}, which
sends any command to any target.

\begin{screen}
10 UCI UCTRL, 1\\
20 PRINT UDAT\$
\end{screen}

That sends command~1 to the control target, which identifies itself, and prints
what came back.

Arguments follow the command, and the wedge sends each number as a single byte
because that is what nearly every command wants. When a command wants something
wider, say so with \bas{uw()} for sixteen bits or \bas{ul()} for thirty-two:

\begin{screen}
10 UCI UDOS1, 16, UW(8192)
\end{screen}

\begin{note}
The SDK knows the shape of many commands and will widen arguments for you
without being told. \bas{uw()} and \bas{ul()} are the escape hatch for a
command newer than the SDK's table, and they always win over it.
\end{note}

\section{What BASIC Cannot Do}

Two limits worth knowing before they surprise you.

\subsection*{IF X THEN UCI does not work}

BASIC's \bas{if} jumps straight into the ROM's own statement dispatcher, which
rejects any token above \bas{new}. Every extension that adds statement keywords
to BASIC~V2 has this hole. So this does not work:

\begin{screen}
10 IF X = 1 THEN ULOAD "GAME"
\end{screen}

and this does:

\begin{screen}
10 IF X = 1 THEN 100\\
\hspace*{1em}\ldots\\
100 ULOAD "GAME"
\end{screen}

Functions are unaffected, because expression evaluation does come back through
the wedge. \bas{if uerr then \ldots} is fine, and so is \bas{if ureu = 0 then
\ldots}.

\subsection*{Strings are 255 bytes}

\bas{udat\$} gives you the reply as a string, clipped at 255 bytes because that
is the longest string BASIC has. When a reply is longer than that, use
\bas{ubyte(n)} to walk it, or \bas{ubload} to put it somewhere you can
\bas{peek}.
